% rho is a distribution over individual neurons; Q is over complete networks.
% Volume three architecture figures. Independent export only; not book fonts.
\colorlet{MutedBlueFill}{FigureBlueFill}
\colorlet{MutedGreenFill}{FigureGreenFill}
\colorlet{MutedPurple}{FigurePurple}
\colorlet{MutedPurpleFill}{FigurePurpleFill}
\tikzset{
 v3 picture/.style={x=1mm,y=1mm,font=\figurefont\fontsize{8}{10}\selectfont,text=FigureInk,line width=.65pt},
 v3 box/.style={draw=FigureBlue,fill=white,rounded corners=1.8pt,line width=.7pt,inner sep=2.5pt,align=center},
 v3 data/.style={v3 box,fill=FigureBlueFill},
 v3 model/.style={v3 box,draw=FigureGreen,fill=FigureGreenFill},
 v3 output/.style={v3 box,draw=FigureYellow,fill=FigureYellowFill},
 v3 risk/.style={v3 box,draw=FigureRed,fill=FigureRedFill},
 v3 group/.style={draw=FigureBlue,rounded corners=1.8pt,line width=.65pt,inner sep=3mm},
 v3 flow/.style={draw=FigureStroke,line width=.75pt,-{Stealth[length=1.65mm,width=1.05mm]},rounded corners=1pt},
 v3 guide/.style={draw=FigureMuted,line width=.6pt,dashed},
 v3 feedback/.style={v3 flow,draw=FigureGreen,dashed},
 v3 title/.style={draw=none,fill=none,font=\figurefont\fontsize{8}{10}\selectfont,align=center,inner sep=1pt},
 v3 note/.style={draw=none,fill=none,font=\FigureNoteFont\fontsize{7.5}{9.5}\selectfont,text=FigureMuted,align=center,inner sep=1pt},
 v3 label/.style={draw=none,fill=none,font=\figurefont\fontsize{8}{10}\selectfont,align=center,inner sep=1pt},
 v3 wire/.style={draw=FigureMuted,line width=.45pt},
 v3 point/.style={circle,draw=FigureBlue,fill=FigureBlueFill,minimum size=2.3mm,inner sep=0pt,line width=.55pt}
}
% Bring pre-existing native TikZ geometry to the same typography and scale.
\tikzset{
 foundation picture/.append style={font=\figurefont\fontsize{8}{10}\selectfont,line width=.65pt},
 foundation node/.append style={rounded corners=1.8pt,line width=.7pt,font=\figurefont\fontsize{8}{10}\selectfont},
 foundation heading/.append style={draw=none,fill=none,font=\figurefont\fontsize{8}{10}\selectfont},
 foundation note/.append style={draw=none,fill=none,font=\FigureNoteFont\fontsize{7.5}{9.5}\selectfont},
 foundation flow/.append style={line width=.75pt,-{Stealth[length=1.65mm,width=1.05mm]}},
 foundation frame/.append style={rounded corners=1.8pt,line width=.7pt},
 foundation edge label/.append style={draw=none,fill=none,font=\figurefont\fontsize{8}{10}\selectfont}
}
% Official Lucide outline icons, consistently used across diagrams. Size in mm.
\newcommand{\VIcon}[5]{\node[inner sep=0pt] at (#2,#3) {\includegraphics[width=#4mm]{figures/03-foundations/00-shared/tikz/icons/#1-#5.pdf}};}
% A neutral network glyph (not a claim about an actual model architecture).
% Coordinates denote its centre; scale 1 gives 13 x 8 mm.
\newcommand{\VNet}[4]{\begin{scope}[shift={(#1,#2)},scale=#3]
 \foreach \a in {-3,0,3}{\foreach \b in {-3,0,3}{\draw[v3 wire] (-5,\a)--(0,\b);}}
 \foreach \a in {-3,0,3}{\draw[v3 wire] (0,\a)--(5,0);}
 \foreach \a in {-3,0,3}{\draw[draw=#4,fill=#4,line width=.55pt] (-5,\a) circle(.9);\draw[draw=#4,fill=white,line width=.55pt] (0,\a) circle(.9);}
 \draw[draw=#4,fill=#4,line width=.55pt] (5,0) circle(.9);
\end{scope}}
\newcommand{\VDots}[4]{\foreach \k in {0,1,2,3}{\fill[#4] ({#1+2.7*\k},#2) circle(#3);}}
% Feature vector/tensor elements are scientific glyphs, not decorative icons.
\newcommand{\VVector}[4]{\foreach \k in {0,...,4}{\fill[#4,rounded corners=.25pt] ({#1+2.2*\k},#2) rectangle ++(1.5,#3);}}

\colorlet{VNetworkInput}{FigureBlue}\colorlet{VNetworkHidden}{FigureYellow}\colorlet{VNetworkOutput}{FigureGreen}
\begin{tikzpicture}[v3 picture]
\tikzset{v3 box/.append style={line width=.8pt},v3 group/.append style={line width=.85pt},v3 flow/.append style={line width=.9pt}}
\path[use as bounding box] (0,0) rectangle (163,79);
\node[v3 title] at (38,74) {均值场：神经元参数分布};
\node[v3 title] at (123,74) {PAC--Bayes：完整网络的后验};
\draw[v3 guide] (79,6)--(79,69);
% Annotation leaders terminate at their actual neuron/operator ports.

% A single network, shown neuron by neuron. Yellow highlights the integration unit.
\node[v3 label] at (8,64) {输入 $x$};
\node[v3 label] at (26,64) {单个神经元};
\node[v3 label] at (54,64) {确定性预测器 $f_\rho$};
\draw[v3 guide,draw=FigureYellow,fill=FigureYellowFill,rounded corners=1.8pt] (19,26) rectangle (29,59);
\foreach \a in {34,44,52}{\foreach \b in {28,34,44,50,56}{\draw[v3 wire] (8,\a)--(24,\b);}}
\foreach \a in {28,34,44,50,56}{\draw[v3 wire] (24,\a)--(37,44);}
\foreach \a in {34,44,52}{\draw[draw=FigureBlue,fill=FigureBlueFill] (8,\a) circle(1.55);}
\foreach \a in {28,34,44,50,56}{\draw[draw=FigureYellow,fill=FigureYellow] (24,\a) circle(1.55);}
\draw[draw=FigureGreen,fill=FigureGreen] (37,44) circle(1.9);
\foreach \y in {37.3,38.4,39.5}{\fill[FigureStroke] (24,\y) circle(.28);}\foreach \y in {37,38.2,39.4}{\fill[FigureStroke] (8,\y) circle(.28);}\node[v3 box,draw=FigureGreen,minimum width=10mm,minimum height=12mm] (int) at (51,44) {$\int$};
\draw[v3 flow,draw=FigureBlue] (8,61)--(8,53.55);\draw[v3 flow,draw=FigureYellow] (24,61)--(24,57.55);
\draw[v3 flow,draw=FigureGreen] (51,61)--(int.north);
% The integral is an alternative representation of the same network, not an ensemble.
\draw[v3 guide] (39,44)--(int.west);
\node[v3 label] at (42.5,47) {$=$};
\node[circle,draw=none,fill=FigureRed,minimum size=3.2mm,inner sep=0pt] (mean-output) at (69,44) {};\draw[v3 flow] (int.east)--(mean-output.west);
\node[v3 label,anchor=north] at (51,35) {平均／积分};\node[v3 label] at (69,55) {输出};
\draw[v3 flow,draw=FigureRed] (69,53)--(mean-output.north);
% Each point below is the parameter of one neuron, never a whole-network sample.
\draw[v3 box,draw=FigureYellow,fill=FigureYellowFill,rounded corners=3mm] (6,10) rectangle (40,23);
\foreach \x/\y in {9/15,11/19,13/12,15/17,18/20,19/13,21/17,23/12,26/19,28/15,30/12,32/18,35/14,37/19}{\draw[draw=FigureYellow,fill=FigureYellow] (\x,\y) circle(.85);}
\draw[v3 flow,draw=FigureYellow] (24,23)--(24,26);
\draw[v3 guide,draw=FigureYellow] (9,23)--(19,26);\draw[v3 guide,draw=FigureYellow] (37,23)--(29,26);
\node[v3 label,anchor=north] at (24,9) {神经元参数分布 $\rho$};
\node[v3 note,anchor=north] at (38,3) {分布元素是单个神经元的参数};
\node[v3 note] at (57,25) {$f_\rho(x)=\int a\phi(w^\top x)\,\mathrm d\rho$};
% Four glyphs are four complete-network hypotheses in the posterior Q.
\node[v3 label] at (104,64) {完整网络的后验 $Q$};
\draw[v3 box,fill=FigureBlueFill] (84,24) rectangle (125,59);
\VLayerNet{94}{50}{.85}{2}{4}{1}\VLayerNet{115}{50}{.85}{2}{4}{1}
\VLayerNet{94}{32}{.85}{2}{4}{1}\VLayerNet{115}{32}{.85}{2}{4}{1}
\node[v3 label] at (105,41) {$\cdots$};
\node[v3 box,draw=FigureGreen,minimum width=19mm,minimum height=16mm] (sample) at (144,41) {};
\VLayerNet{144}{41}{1}{2}{4}{1}
\node[v3 label] at (144,60) {抽取的网络};
\draw[v3 flow,dashed] (125,41)--(sample.west);
\node[v3 label,anchor=north] at (133,29) {抽取 $\theta\sim Q$};
\node[circle,draw=none,fill=FigureRed,minimum size=3.2mm,inner sep=0pt] (sample-output) at (160,41) {};\draw[v3 flow] (sample.east)--(sample-output.west);
\node[v3 label] at (159,54) {输出};\draw[v3 flow,draw=FigureRed] (159,52)--(sample-output);
\node[v3 note] at (124,13) {分布元素是完整网络的参数};
\node[v3 label] at (124,5) {$\theta\sim Q\quad\longrightarrow\quad f_\theta(x)$};
\end{tikzpicture}
